% ============================================================================
% Master Techniques (Test-Taking Tips) — 7 Practice Tests: The Forge
% Blacksmith techniques with the STRIKE method and tactical approach
% Grade 7
% ============================================================================

\thispagestyle{empty}

{
\sffamily

% --- Title ---
\begin{center}
    \textcolor{funOrangeDark}{\fontsize{26}{30}\selectfont\bfseries
    \faShieldAlt~Master Techniques~\faShieldAlt}

    \smallskip
    \textcolor{funGray}{\normalsize The techniques every skilled crafter needs at the anvil}
\end{center}

\vspace{3mm}

% --- The STRIKE Method (unique to this book) ---
\begin{tcolorbox}[
    enhanced,
    colback=funGrayDark!8,
    colframe=funOrangeDark!60,
    arc=10pt, boxrule=2pt,
    left=5mm, right=5mm, top=4mm, bottom=4mm,
    title={\textcolor{white}{\large\faHammer~The S.T.R.I.K.E.\ Method}},
    colbacktitle=funGrayDark!80,
    fonttitle=\bfseries\sffamily,
]
\small
\renewcommand{\arraystretch}{1.5}
\begin{tabular}{@{}c@{\hspace{3mm}}l@{\hspace{3mm}}p{9.5cm}@{}}
    \textcolor{funOrange}{\fontsize{18}{18}\selectfont\bfseries S} &
    \textbf{Scan} &
    Read the whole problem. What is given? What is asked? \\
    \textcolor{funOrange}{\fontsize{18}{18}\selectfont\bfseries T} &
    \textbf{Tag} &
    Underline key values. Circle the operation clue words. \\
    \textcolor{funOrange}{\fontsize{18}{18}\selectfont\bfseries R} &
    \textbf{Reason} &
    Choose a strategy: proportion, equation, diagram, or estimation. \\
    \textcolor{funOrange}{\fontsize{18}{18}\selectfont\bfseries I} &
    \textbf{Implement} &
    Solve step by step. Write every calculation. Label units. \\
    \textcolor{funOrange}{\fontsize{18}{18}\selectfont\bfseries K} &
    \textbf{Know} &
    Check: does the answer make sense in context? \\
    \textcolor{funOrange}{\fontsize{18}{18}\selectfont\bfseries E} &
    \textbf{Evaluate} &
    Re-read the question. Did you answer what was actually asked? \\
\end{tabular}
\end{tcolorbox}

\vspace{2mm}

% --- Forging Tips (3 columns of technique cards) ---
\noindent\textcolor{funOrangeDark}{\large\bfseries\faFire~Forging Techniques}

\vspace{2mm}

\noindent
\begin{minipage}[t]{0.48\linewidth}
    \begin{tcolorbox}[
        enhanced,
        colback=funRed!5,
        colframe=funRed!30,
        arc=6pt, boxrule=1pt,
        left=3mm, right=3mm, top=2mm, bottom=2mm,
    ]
        \small
        \textcolor{funRedDark}{\bfseries Technique 1:}
        Convert between fractions, decimals, and percents
        freely. Many Grade 7 problems become easier in a different form.
    \end{tcolorbox}

    \vspace{1.5mm}

    \begin{tcolorbox}[
        enhanced,
        colback=funOrange!5,
        colframe=funOrange!30,
        arc=6pt, boxrule=1pt,
        left=3mm, right=3mm, top=2mm, bottom=2mm,
    ]
        \small
        \textcolor{funOrangeDark}{\bfseries Technique 2:}
        When solving proportions, cross-multiply carefully.
        Keep track of which quantity is in the numerator and
        which is in the denominator.
    \end{tcolorbox}

    \vspace{1.5mm}

    \begin{tcolorbox}[
        enhanced,
        colback=funYellow!8,
        colframe=funYellowDark!35,
        arc=6pt, boxrule=1pt,
        left=3mm, right=3mm, top=2mm, bottom=2mm,
    ]
        \small
        \textcolor{funYellowDark}{\bfseries Technique 3:}
        For integer operations, use the number line mentally.
        Subtracting a negative? That's adding. Don't let
        sign rules trip you up.
    \end{tcolorbox}
\end{minipage}%
\hfill
\begin{minipage}[t]{0.48\linewidth}
    \begin{tcolorbox}[
        enhanced,
        colback=funGreen!5,
        colframe=funGreen!30,
        arc=6pt, boxrule=1pt,
        left=3mm, right=3mm, top=2mm, bottom=2mm,
    ]
        \small
        \textcolor{funGreenDark}{\bfseries Technique 4:}
        On geometry questions, draw a diagram if one isn't given.
        Label all known measurements before computing.
    \end{tcolorbox}

    \vspace{1.5mm}

    \begin{tcolorbox}[
        enhanced,
        colback=funBlue!5,
        colframe=funBlue!30,
        arc=6pt, boxrule=1pt,
        left=3mm, right=3mm, top=2mm, bottom=2mm,
    ]
        \small
        \textcolor{funBlueDark}{\bfseries Technique 5:}
        Estimation is your ally. Before computing an exact answer,
        estimate. If your answer is far from the estimate,
        recheck your work.
    \end{tcolorbox}

    \vspace{1.5mm}

    \begin{tcolorbox}[
        enhanced,
        colback=funPurple!5,
        colframe=funPurple!30,
        arc=6pt, boxrule=1pt,
        left=3mm, right=3mm, top=2mm, bottom=2mm,
    ]
        \small
        \textcolor{funPurpleDark}{\bfseries Technique 6:}
        Never leave a blank. On constructed response,
        even writing the equation you'd use can earn
        partial credit.
    \end{tcolorbox}
\end{minipage}

\vspace{3mm}

% --- Pre-forge ritual ---
\begin{tcolorbox}[
    enhanced,
    colback=funOrange!4,
    colframe=funOrange!20,
    arc=8pt, boxrule=0.8pt,
    left=5mm, right=5mm, top=3mm, bottom=3mm,
]
    \small\sffamily
    \textcolor{funOrangeDark}{\bfseries\faMoon~Night Before Practice:}
    Sleep 8+ hours.
    Eat a solid meal.
    Gather your tools (pencils, eraser, scratch paper).
    Tell yourself: \textit{``Every strike makes me stronger.''}
\end{tcolorbox}

\vspace{2mm}

% --- Forgemaster encouragement ---
\begin{tcolorbox}[
    enhanced,
    colback=funYellow!8,
    colframe=funYellowDark!25,
    arc=10pt, boxrule=0.8pt,
    left=5mm, right=5mm, top=3mm, bottom=3mm,
]
    \begin{minipage}[c]{0.07\linewidth}
        \centering \iconOwl[1cm]
    \end{minipage}%
    \hfill
    \begin{minipage}[c]{0.90\linewidth}
        \small\itshape\color{funOrangeDark}
        \textbf{The Forgemaster says:} ``The greatest crafters don't avoid
        the hardest pieces --- they lean into them. Review every mistake.
        That's where the real tempering happens.''
    \end{minipage}
\end{tcolorbox}

}

\newpage
